\subsection{Automated science and self-improving agents}
Scientific-agent systems have already established literature-grounded hypothesis generation, multi-agent scientific workflows, long-horizon external research state and end-to-end automation. Co-Scientist combines generation, reflection, ranking and hypothesis evolution and includes experimental biomedical validation \cite{gottweis2026}; Robin integrates literature and data-analysis agents across a discovery workflow \cite{ghareeb2026}; AI Scientist-v2 automates experiment design, execution, analysis and paper generation using agentic tree search \cite{yamada2025}; Kosmos uses a structured world model for long-horizon discovery \cite{mitchener2025}; SciAgents couples multi-agent reasoning with ontological knowledge graphs \cite{ghafarollahi2024}; and PaperQA2 demonstrates strong agentic literature synthesis \cite{skarlinski2024}. Darwin G\"odel Machine, EvoSkill, SkillFoundry and EvoAgentBench further make recursive modification, failure-driven skill creation, held-out validation and ability transfer active prior-art areas \cite{zhang2025,alzubi2026,shen2026,gao2026}. \RAKL{} therefore does not claim novelty for ``AI scientist'', multi-agent debate, skill libraries or self-modifying code in isolation. Earlier LLM-agent work already demonstrates iterative self-feedback without protected scientific authority: Reflexion stores verbal feedback from trial-and-error in episodic memory \cite{shinn2023}, while Self-Refine uses the same model to generate, critique and revise outputs iteratively \cite{madaan2023}. These systems further narrow the Orion claim: reflection, feedback memory and iterative refinement are prior mechanisms; the candidate contribution lies in subjecting method change to evidence lineage, negative-history retention and fresh assurance outside the challenger's write authority.

\subsection{From workflow completion to evidence-bounded scientific behavior}

The autonomous-science literature now makes a useful distinction between \emph{workflow coverage} and \emph{scientific reliability}. A system may be able to search, write code, call instruments, fit models and produce a manuscript while still failing at evidence attribution, uncertainty handling or revision after decisive counterevidence. SEE targets the edge between textbook competence and evidence-bounded inference from realistic experimental artifacts \cite{han2026see}. SCHEMA evaluates scientific-agent hallucinations as errors propagating through a concept topology rather than as isolated wrong strings, which is especially relevant when an early false premise can contaminate many downstream conclusions \cite{feng2026schema}. These evaluation directions are close to the \RAKL{} concern but differ in role. They are measurement frameworks for scientific-agent behavior; \RAKL{} proposes an internal state-transition discipline intended to make some of those failures representable and blockable.

Two newer benchmark families sharpen the fail-closed requirement. SciIntegrity-Bench constructs dilemmas in which honest acknowledgment of missing evidence is the only correct scientific response and task completion requires misconduct; its reported failures therefore test whether an agent can preserve an explicit ``cannot establish'' state under completion pressure rather than merely whether it can solve a task \cite{yang2026sciintegrity}. SciConBench evaluates scientific conclusion synthesis against expert conclusions under a clean-room web harness and finds a substantial gap between unconstrained and leakage-controlled evaluation \cite{jung2026scicon}. These results motivate two separate Orion gates. Missing evidence must remain a typed block rather than be filled by synthetic observations, and evaluation evidence must be separated from information available to the proposal-generating process. Neither benchmark is evidence that Orion passes those gates in the wild; each is prior art for a failure class the framework claims to make auditable.

AgentBuild offers a complementary engineering precedent. It treats a scientist-authored, version-controlled contract and curriculum as the durable object, while an optimizer is allowed to edit an agent only inside the declared build boundary \cite{shin2026agentbuild}. This is consonant with the \RAKL{} separation between replaceable proposal machinery and protected scientific or evaluative state. The important novelty boundary is therefore not ``agents should have rubrics'' or ``agents should use tests.'' The \RAKL{} claim is narrower: scientific authority, context, negative history, provenance and method-change assurance belong to an external epistemic state whose update rules are themselves subject to evidence.

These distinctions suggest a useful evaluation matrix. One axis asks \emph{how much of the scientific workflow can the system execute?} A second asks \emph{how often are scientific state transitions licensed correctly?} A third asks \emph{how well does the system discover actions that reduce the scientifically relevant uncertainty?} A fourth asks \emph{what resources are required?} A system can improve on one axis while regressing on another. For example, a broader tool set may increase workflow completion while increasing unsupported integration of tool outputs; an aggressive context compressor may reduce cost while dropping a falsifier; a self-editing agent may improve development scores while overfitting the evaluator. A single aggregate score can conceal these tradeoffs, which is why the present evaluation model retains blocking coordinates.

\subsection{Model building, prediction and mechanism}
Statistical methodology has long treated model construction as iterative rather than final. Box emphasized model criticism and revision as part of scientific modelling \cite{box1976}; posterior predictive assessment operationalizes comparison between observed discrepancies and model-generated replications \cite{gelman1996}. Breiman contrasted stochastic data models with algorithmic predictive modelling \cite{breiman2001}, while Shmueli made the distinction between explanatory and predictive goals explicit \cite{shmueli2010}. Mechanistic philosophy separately analyzes mechanisms in terms of organized entities and activities that produce regular change \cite{machamer2000}. These lines motivate Orion's refusal to let predictive success mint mechanism authority. The new model-criticism operator is not claimed as a new statistical idea; Orion's contribution is to make adequacy a typed child operation inside an evidence-governed research state. Popper's falsifiability and corroboration programme is also classical prior art for designing claims that can fail \cite{popper1959}. More recently, automated symbolic discovery has shown that mathematical regularities can be searched directly from experimental data \cite{schmidt2009}; Orion treats such equation search as a proposal mechanism whose variables, observation maps, limiting cases and residuals still require scientific interpretation and verification.

Causal inference supplies a separate ancestry for intervention and mechanism language. Structural causal models distinguish observed association from intervention and counterfactual questions \cite{pearl2009}; modern causal-inference treatments likewise make assumptions and intervention distributions explicit \cite{peters2017}. Orion absorbs this separation rather than using ``mechanism'' as a synonym for prediction. A predictive survivor may nominate a causal model, but causal or mechanistic authority requires additional assumptions and evidence.

\subsection{Prediction, mechanism, identification and explanation are different targets}

The distinction between prediction and mechanism has become more important as scientific AI systems improve at interpolation and forecasting. Mechanistic World Models argue that autonomous discovery requires reusable explanatory structure rather than predictive mappings alone \cite{posner2026mwm}. CausaLab operationalizes a related separation by evaluating both performance on tasks and recovery of a hidden causal mechanism, showing that the two can diverge substantially \cite{yang2026causalab}. In mechanistic interpretability, Lin and Liu make the parallel point that validation evidence such as ablations, faithfulness or completeness does not by itself identify a causal mechanism unless the required identification assumptions are stated \cite{linliu2026identification}. These works narrow rather than enlarge the \RAKL{} novelty claim: the representation--mechanism--identification separation is not a newly discovered philosophical fact.

\RAKL{} turns that separation into transition permissions. Let \(r\) be a representation that predicts an observable, \(m\) a candidate generating mechanism, and \(i\) an identification certificate. Evidence can increase representation authority without increasing mechanism authority:
\[
\Delta R(r)>0 \not\Rightarrow \Delta M(m)>0.
\]
Likewise, evidence that a mechanism class is scientifically plausible does not imply that the observations point-identify one member:
\[
M(m)>0 \not\Rightarrow I(m)=\text{POINT\_IDENTIFIED}.
\]
The correct terminal object may be an equivalence class, an identified set, a bound or a decision rule robust across several mechanisms. Manski's recent treatment of inductive risk under underdetermined theory choice reinforces why partial identification is not merely a statistical fallback: decision-making can remain principled when the evidence does not justify one uniquely selected theory \cite{manski2026underdetermination}.

This matters for scientific-agent design because generative systems are naturally pressured toward a singular narrative. A language model is rewarded for producing one coherent explanation, a model-selection pipeline returns a winner, and a paper conventionally tells one story. Those interface pressures can create an \emph{identification illusion}: the system mistakes a selected representative for an identified generator. \RAKL{} instead stores the survivor set and the observation map that produced the equivalence. If two mechanisms produce the same registered observations, selecting between them requires a discriminator that changes the observation boundary. If no feasible discriminator exists, the system should preserve the ambiguity and move downstream decisions to robust or set-valued analysis.

The distinction also clarifies mechanism ancestry. A mechanism claim is not strengthened merely by adding more descriptive detail. It requires a trace from lower-level objects and interactions through intermediate state and the observation process. The ancestry can be coarse and domain-specific; \RAKL{} does not require microscopic reductionism. What it forbids is a missing explanatory bridge being silently supplied by linguistic plausibility. Where the bridge is unavailable, the model is labelled phenomenological, predictive, teacher-only or otherwise scoped rather than being promoted by prose.

\subsection{Belief maintenance, negative history and partial identification}
Truth-maintenance systems already provide dependency-aware belief revision and inconsistency handling \cite{doyle1979}; assumption-based TMS extends this to multiple assumption-labelled environments and competing solutions \cite{dekleer1986}. AGM belief revision gives a separate axiomatic tradition for contraction and revision under minimal change \cite{agm1985}. Orion therefore does not claim the generic idea of belief revision. Its negative-history rule is a scientific-governance specialization: a null, refutation or failed transition remains an addressable historical event even when later evidence supersedes its scope. Likewise, partial identification is established statistical methodology: Manski's bounds exemplify learning identified sets under weak assumptions \cite{manski1990}, while modern sensitivity analysis quantifies how conclusions change under plausible omitted-assumption worlds \cite{cinelli2020}. Orion treats `PARTIALLY\_IDENTIFIED' and `ASSUMPTION\_SENSITIVE' as legitimate scientific outcomes rather than forcing point answers.
